\documentclass[12pt]{article}
\usepackage{setspace}
\usepackage{amsmath,amssymb,euscript,graphicx}
\usepackage{xurl}
\usepackage{graphicx}
\usepackage{subfigure}
\usepackage{wrapfig}
\usepackage[compact]{titlesec}
\usepackage{fancybox,color}
\usepackage[footnotesize]{caption}
\usepackage{cite}
\usepackage{enumitem}
\usepackage[normalem]{ulem}
\usepackage[top=0.9375in, right=0.9375in, left=0.9375in, bottom=0.9375in, centering]{geometry}
\usepackage{bm}
\usepackage{float}

\usepackage{amsthm}


% -- Loading the code block package:
\usepackage{listings}
% -- Basic formatting
\usepackage[utf8]{inputenc}
\usepackage[english]{babel}
\usepackage{times}
\setlength{\parindent}{8pt}
\usepackage{indentfirst}
% -- Defining colors:
\usepackage[dvipsnames]{xcolor}
\definecolor{codegreen}{rgb}{0,0.6,0}
\definecolor{codegray}{rgb}{0.5,0.5,0.5}
\definecolor{codepurple}{rgb}{0.58,0,0.82}
\definecolor{backcolour}{rgb}{0.95,0.95,0.92}
% Definig a custom style:
\lstdefinestyle{mystyle}{
    backgroundcolor=\color{backcolour},   
    commentstyle=\color{codepurple},
    keywordstyle=\color{NavyBlue},
    numberstyle=\tiny\color{codegray},
    stringstyle=\color{codepurple},
    basicstyle=\ttfamily\footnotesize\bfseries,
    breakatwhitespace=false,         
    breaklines=true,                 
    captionpos=t,                    
    keepspaces=true,                 
    numbers=left,                    
    numbersep=5pt,                  
    showspaces=false,                
    showstringspaces=false,
    showtabs=false,                  
    tabsize=2
}
% -- Setting up the custom style:
\lstset{style=mystyle}


% IEEE uses Times Roman font, so we'll default to Times.
% These three commands make up the entire times.sty package.
\renewcommand{\sfdefault}{phv}
\renewcommand{\rmdefault}{ptm}
\renewcommand{\ttdefault}{pcr}

\newcommand{\squeeze}{\vspace{-0.1in}}
\newcommand{\changed}[1]{\textcolor{blue}{#1}} %generates text in blue

%%%%%%%%%%%%%%%%%%%%%%%%%%%%%%%%%%%%%%%%%%%%%%%%%%%%%%%%%%%%%%%%%%%%%%%

\begin{document}

\begin{center}
  {\Large {\bf Lab 3: Making use of for loops}}
  
\end{center}

\vspace{5pt}
\hrule
\vspace{5pt}

\section{Summary}
In this lab, we will reduce the number of lines of code by taking advantage of for loops. 

\section{Setup}
\begin{itemize}
  \item open a terminal in VS Code. 
  \item Navigate to your ES1098 folder, use \texttt{ls}, \texttt{cd}, and \texttt{pwd} to help you find this directory
  \item Once you are in this directory you should see: \texttt{/Users/<computername>/ES1098} when you type \texttt{pwd}.
  \item Create a new folder by typing \texttt{mkdir lab3}
  \item Move into that folder by typing \texttt{cd lab3}
  \item Confirm your current directory by typing \texttt{pwd}.
\end{itemize}

In VS Code, make a new file called \texttt{lab2.py}. Remember: you can do this using the \texttt{code lab3.py} in the terminal of your VS Code IDE. Make sure you are in the correct directory. 

\section{Using Loops to define simple shapes}

\section{Design a repeating Geometric pattern}



\section{Acknowledgments}
This lab is from Professor Hannen Wolfe in the Colby College CS Department. 
Sources:
\begin{itemize}
\item \footnotesize{\url{https://cs.colby.edu/hewolfe/courses/cs151/f21/projects/02/Lab02FunctionsToControlShapes.html}}
\end{itemize}


\end{document}
